%        File: revise.tex
%     Created: Wed Oct 27 02:00 PM 2018 P
% Last Change: Wed Oct 27 02:00 PM 2018 P
%

%
% Copyright 2007, 2008, 2009 Elsevier Ltd
%
% This file is part of the 'Elsarticle Bundle'.
% ---------------------------------------------
%
% It may be distributed under the conditions of the LaTeX Project Public
% License, either version 1.2 of this license or (at your option) any
% later version.  The latest version of this license is in
%    http://www.latex-project.org/lppl.txt
% and version 1.2 or later is part of all distributions of LaTeX
% version 1999/12/01 or later.
%
% The list of all files belonging to the 'Elsarticle Bundle' is
% given in the file `manifest.txt'.
%

% Template article for Elsevier's document class `elsarticle'
% with numbered style bibliographic references
% SP 2008/03/01
%
%
%
% $Id: elsarticle-template-num.tex 4 2009-10-24 08:22:58Z rishi $
%
%
%\documentclass[preprint,12pt]{elsarticle}
\documentclass[answers,11pt]{exam}

% \documentclass[preprint,review,12pt]{elsarticle}

% Use the options 1p,twocolumn; 3p; 3p,twocolumn; 5p; or 5p,twocolumn
% for a journal layout:
% \documentclass[final,1p,times]{elsarticle}
% \documentclass[final,1p,times,twocolumn]{elsarticle}
% \documentclass[final,3p,times]{elsarticle}
% \documentclass[final,3p,times,twocolumn]{elsarticle}
% \documentclass[final,5p,times]{elsarticle}
% \documentclass[final,5p,times,twocolumn]{elsarticle}

% if you use PostScript figures in your article
% use the graphics package for simple commands
% \usepackage{graphics}
% or use the graphicx package for more complicated commands
\usepackage{graphicx}
% or use the epsfig package if you prefer to use the old commands
% \usepackage{epsfig}

% The amssymb package provides various useful mathematical symbols
\usepackage{amssymb}
% The amsthm package provides extended theorem environments
% \usepackage{amsthm}
\usepackage{amsmath}

% The lineno packages adds line numbers. Start line numbering with
% \begin{linenumbers}, end it with \end{linenumbers}. Or switch it on
% for the whole article with \linenumbers after \end{frontmatter}.
\usepackage{lineno}

% I like to be in control
\usepackage{placeins}

% natbib.sty is loaded by default. However, natbib options can be
% provided with \biboptions{...} command. Following options are
% valid:

%   round  -  round parentheses are used (default)
%   square -  square brackets are used   [option]
%   curly  -  curly braces are used      {option}
%   angle  -  angle brackets are used    <option>
%   semicolon  -  multiple citations separated by semi-colon
%   colon  - same as semicolon, an earlier confusion
%   comma  -  separated by comma
%   numbers-  selects numerical citations
%   super  -  numerical citations as superscripts
%   sort   -  sorts multiple citations according to order in ref. list
%   sort&compress   -  like sort, but also compresses numerical citations
%   compress - compresses without sorting
%
% \biboptions{comma,round}

% \biboptions{}


\usepackage{xspace}
\usepackage{color}

\usepackage{multirow}
\usepackage[hyphens]{url}


\usepackage[acronym,toc]{glossaries}
\include{acros}

\makeglossaries

%\journal{Annals of Nuclear Energy}

\begin{document}

%\begin{frontmatter}

% Title, authors and addresses

% use the tnoteref command within \title for footnotes;
% use the tnotetext command for the associated footnote;
% use the fnref command within \author or \address for footnotes;
% use the fntext command for the associated footnote;
% use the corref command within \author for corresponding author footnotes;
% use the cortext command for the associated footnote;
% use the ead command for the email address,
% and the form \ead[url] for the home page:
%
% \title{Title\tnoteref{label1}}
% \tnotetext[label1]{}
% \author{Name\corref{cor1}\fnref{label2}}
% \ead{email address}
% \ead[url]{home page}
% \fntext[label2]{}
% \cortext[cor1]{}
% \address{Address\fnref{label3}}
% \fntext[label3]{}

\title{Hybrid $S_N$-Diffusion Neutronics Method for Control Rod Modeling in
        Molten Salt Reactors\\
        \large Response to Review Comments}
\author{Sun Myung Park, Kathryn D.\ Huff, Madicken Munk}
\date{}

% use optional labels to link authors explicitly to addresses:
% \author[label1,label2]{<author name>}
% \address[label1]{<address>}
% \address[label2]{<address>}


%\author[uiuc]{Kathryn Huff}
%        \ead{kdhuff@illinois.edu}
%  \address[uiuc]{Department of Nuclear, Plasma, and Radiological Engineering,
%        118 Talbot Laboratory, MC 234, Universicy of Illinois at
%        Urbana-Champaign, Urbana, IL 61801}
%
% \end{frontmatter}
\maketitle
\section*{Review General Response}
We sincerely thank the reviewers for their careful evaluation of our manuscript and for the
detailed and constructive feedback. The comments significantly improved the clarity, technical
rigor, and positioning of the work. We have revised the manuscript accordingly and address each
comment in this response document.

Major revisions include:
%
\begin{enumerate}
  \item Clarification of Scope and Claims

We have revised the Introduction, Abstract, and Conclusion to more clearly state that the hybrid
method is not intended to recover full-domain transport accuracy. Instead, it provides localized
transport correction near control rods for improved rod worth and power prediction in multiphysics
simulations. Claims regarding efficiency and accuracy have been moderated accordingly.

  \item Improved Error Attribution

    The contribution of vacuum boundary leakage to $k_\text{eff}$ discrepancies has been quantified and
explicitly discussed in the Results section. We clarify that the hybrid method does not correct
leakage effects on the outer boundary.
We also now clearly separate the effects of energy discretization, spatial discretization, and
leakage when interpreting the results.

  \item Clarifications on Method Implementation

The descriptions of outer iterations, convergence criteria, adaptive buffer region boundary, and the
diffusion coefficient limiter have been revised to remove ambiguity.

  \item Reduced Overstatement of Efficiency

Computational cost comparisons have been reframed to emphasize that the hybrid method incurs
moderate additional cost relative to diffusion while avoiding the significantly larger cost of
full-domain $S_8$ transport.
\end{enumerate}

We believe these revisions substantially strengthen the manuscript and address all reviewer concerns.

%---------------------------------------------------------------------
% REVIEWER 2
%---------------------------------------------------------------------
\section*{Reviewer 2}
\begin{questions}

    %--- Abstract Comments ---
    \question (Abstract Comment) Minor comment, but I would reverse and tweak the first two sentences to better set the overall stage.
    \begin{solution}
        Thank you for your suggestion. We edited the two sentences as recommended.
    \end{solution}

    \question (Abstract Comment) You might reconsider using the word "novel." Generations of people have done various diffusion/transport hybrid methods, for various physics applications, using diffusion where they can get away with it and transport where they have to.
    \begin{solution}
        We removed the word ``novel''.
    \end{solution}

    \question (Abstract Comment) The fact that you used MOOSE as the framework for doing this should be touted, I would think. It would be interesting if you could relay some estimate of the cost to develop this capability and compare it to other known efforts. Did it save time in development, testing, validating, etc.?
    \begin{solution}
        We added the following sentence to the abstract: ``We implemented the hybrid method in
        Moltres, a MOOSE-based multiphysics reactor analysis application, leveraging MOOSE's robust
        finite-element solver infrastructure together with Moltres' existing reactor modeling
        capabilities to reduce development time.'' We also added text to the introduction and
        numerical implementation sections highlighting the benefits of using MOOSE.
    \end{solution}

    \question (Abstract Comment) The method is accurate. Is it? Most hybrid methods are designed to converge to the transport solution, but yours does not seem to. Why? The HOLO method had some issues in that the consistency terms were so expensive and difficult to calculate. Is that happening here?
    \begin{solution}
        We accept that the accuracy statement may be misleading and so we replaced
        the sentence with:
        ``In 1-D and 2-D $k$-eigenvalue simulations, the multigroup hybrid method reduced control
        rod worth errors relative to multigroup diffusion and improved channel power predictions.''

        If we focus solely on the $S_N$ solution in the subregion, it converges to the equivalent
        $S_N$ transport solution given the same subregion boundary sources as defined by Eq.\ 19
        and 31.

        This hybrid method does not converge to the transport solution in the entire problem domain
        because the transport
        corrections are limited only to the subregion around the control rods. Discrepancies
        arising from uncorrected reaction rates and neutron leakage in the rest of the problem
        domain still exist. We plan to address those discrepancies through more conventional
        approaches such as albedo boundary conditions and the $SP_N$ method since these
        existing approaches are computationally cheaper than the hybrid method.

        Neutron leakage error alone accounts for about -300~pcm
        discrepancy between the hybrid method and OpenMC-MG (multigroup) for the 1-D cases with vacuum
        boundary conditions. Eliminating this source of error would bring our hybrid method keff
        estimates to within $<$200~pcm range of OpenMC-MG, which would be reasonable for
        deterministic neutronics codes used in transient multiphysics reactor analyses. Rod worth
        estimates fall within 1\% of OpenMC-MG while the root-mean-square and maximum errors in
        channel power distributions fall within 0.5\% and 3\%. While not perfect, these represent
        significant improvements over the diffusion method.

        We edited text throughout the paper to more clearly
        reflect the discussion points made here.
    \end{solution}

    \question (Abstract Comment) The method is computationally efficient. I did not see one mention of speedups, or any runtimes or anything related to computational cost and improvements.
    \begin{solution}
        We are unable to provide definitive comparisons between the hybrid and $S_N$ methods because the 1-D
        cases were very fast ($<$2 seconds) on single processor while 2-D full-core $S_8$
        calculations required too much memory. We are currently looking into fixes for the excessive
        memory usage by the $S_N$ solver.

        However, we are able to provide comparisons between the hybrid method and the diffusion
        method for 2-D full-core calculations.
        The hybrid method took approximately 4.9 times as long as the diffusion method. This should
        be representative of 3-D calculations because the $S_N$ subsolver did not take advantage of
        symmetries in the 2-D angular domain. We added this observation and discussion to the
        results \& discussion section of the paper.
    \end{solution}

    \question (Abstract Comment) Adaptive boundaries. How is it adaptive? Don't you just manually set the location of the boundary? I may have missed something.
    \begin{solution}
        The physical boundaries of the correction subregion $V_1$
        are indeed manually fixed. However, the ``adaptive'' aspect refers to the buffer region. The
        algorithm dynamically determines the inner boundary of this buffer region during each
        iteration by locating where the drift correction parameter changes sign (i.e., reaches
        zero). We have extensively revised Section II.C to clarify this distinction.
    \end{solution}

    \question (Abstract Comment) In eigenvalue problems, keff, as an integral quantity, converges faster than, say, the eigenmode. I don't know, but it seems like control rod worth estimates might converge even faster. If so, then it is much less constraining to say that you match control rod worths than keff, which is less constraining than saying that you match the flux.
    \begin{solution}
        We recognize that the keff and rod worths converge faster than the spatial flux
        distribution. We highlighted rod worth comparisons in the abstract because we are
        developing this hybrid method for use in multiphysics MSR transient analyses involving
        control rod motion. We added that the hybrid method also improves channel power distribution
        to the abstract.
    \end{solution}

    %--- General Comments ---
    \question Page 15, lines 19-29: I doubt that putting a limiter on a diffusion coefficient actually increases the rate of convergence. It may improve, or enable, stability.
    \begin{solution}
        Thank you for catching this error. We agree that it is more precise to say the limiter
        improves stability. We removed the statement about the rate of convergence and focused on
        the improvement in stability.
    \end{solution}

    \question If your hybrid method is run on a problem where diffusion is perfectly adequate, do you recover the diffusion solution? It doesn't look like those results are in Figure 6, and they should be. That would be very concerning if the hybrid doesn't properly perform a no-op.
    \begin{solution}
        Case 2a, featuring a homogenized fuel-graphite region sandwiched by air and reflector
        regions, shows that we can recover the diffusion solution when using the hybrid
        method on problems where diffusion is adequate. The keff estimates from the hybrid and
        diffusion methods are consistent with each other. Both methods deviate from OpenMC-MG and
        $S_8$ due to neutron leakage errors on the vacuum boundary.
    \end{solution}

    \question No explicit mention of consistency in the hybrid method equations. This is important for recovering the transport solution and maintaining the diffusion solution.
    \begin{solution}
        Please refer to our response to Comment 4.
    \end{solution}

    \question You start out Section II.D saying that the hybrid method requires appropriate diffusion-SN boundary conditions to converge. This is a strong statement and should be backed up with more information. What level of inappropriateness causes the method to not converge? Do you have a mathematical proof? Anecdotal evidence?
    \begin{solution}
        We misused the word ``convergence'' here. The
        original intent of the opening sentence is a simple statement that the $S_N$
        subproblem requires boundary sources. We edited
        the sentence to:
        ``It also requires boundary sources along $\partial V_1$ derived from
        the neutron diffusion solution.''
    \end{solution}

    \question Figure 1 should be clarified. What does "Final ... solutions obtained" in the bottom oval mean? k wasn't mentioned above in the figure at all. Are more calculations made after convergence, or is the last iteration merely retained?
    \begin{solution}
        We mistakenly omitted $k$.
        We overhauled iteration algorithm steps to provide a clearer description.
    \end{solution}

    \question As I was reading the subsolver boundary conditions section, my expectation was that you have to place the boundary in a location where diffusion is valid, but if it is place to far into the diffusion region, you lose efficiency (but not accuracy). There was a lot of discussion about setting the buffer region, but it seems that you're simply looking for an extrapolation distance (such as used in a Marshak boundary condition, which is sometimes transport-corrected from a coefficient of 2/3 to 0.7104). If this is the case, can you make this section more concise and more mathematically rigorous?
    \begin{solution}
        Your analysis is correct; the buffer region effectively acts as an extrapolation distance
        where transport-derived drift coefficients are discarded. We heavily revised Section
        II.C to make it more concise and mathematically rigorous in response to this comment and
        Comment 6.
    \end{solution}

    \question Regarding Figure 3, I was expecting the subsolver region to extend just beyond the control rods to capture any control-rod-related transport effects that extend outward from the control rod. What does it mean that you do, or have to, go out far enough to include two fuel regions? It seems like there are transport effects at every fuel-graphite interface, which would mean that the SN subregion should be centered around the control rod and every fuel-graphite interface. I am not finding that seeing only Group 1 and 8 fluxes is enough. Can you also plot the total flux? It is not clear what motivated the selection of the buffer location. The reference D's in Figure 3 suggest that you shouldn't have stopped at 10cm, but rather set an SN subsolver around every interface.
    \begin{solution}
        While it is true that transport effects exist at every fuel-graphite interface, placing an
        $S_N$ subsolver around every interface would nearly equate to running a full-domain transport
        simulation. We extend the subregion to 10~cm to ensure at least one neutron mean free path
        between the control rod+air region and the subregion boundary. This helped separate
        transport effects arising from the control rod from the $P_1$ boundary effects. The goal is
        to accurately capture the severe angular gradients near the rod to improve rod worth
        estimates, rather than perfectly resolving the entire heterogeneous core.
        We added phrasing in Section II.C to explicitly state our reasoning for the subregion size.
    \end{solution}

    \question At the end of Section II.F, you say, "...for the possibility of overall convergence." That is a very loaded term. Can you give more rigorous details of the convergence or lack of convergence?
    \begin{solution}
        We overhauled iteration algorithm steps and the subsequent paragraph in Section II.C to
        provide clearer descriptions of the convergence criteria.
    \end{solution}

    \question In Table II, under the SN Transport column, can you -- or should you -- be more specific that "discrete" for the first two rows? Like "FE" or "collocation"?
    \begin{solution}
        We replaced all instances of the word ``discrete'' in the table with ``FE'',
        ``Collocation'', or ``Multigroup''.
    \end{solution}

    \question Figure 4 shows a nice progression of test problems. The 1D is slab or Cartesian, correct?
    \begin{solution}
        Thank you. Yes these are 1-D slab problems. We added the ``slab'' label to the revised draft.
    \end{solution}

    \question Figure 5: are the cells still in the diffusion limit as you refine the mesh?
    \begin{solution}
        All mesh cells at all levels of refinement are smaller than the mean free paths in all
        regions.
    \end{solution}

    \question On page 25, line 48, you attributed reactivity discrepancies to the eight-group structure. How? Can you say more? Does it matter? Can you plot the cross-sections, comparing CE with 8-group? What would such a plot tell you? There doesn't seem to be any MG effect for cases 2b and 3b in Figure 6, except for the hybrid method. This is concerning and indicates that MG is causing some sort of problem with the hybrid method.
    \begin{solution}
        We clarified our reasoning for attributing reactivity discrepancies to the eight-group
        structure by comparing OpenMC-CE (continuous) and OpenMC-MG (multigroup) results. OpenMC-MG
        is inconsistent with OpenMC-CE in Cases 2a and 3a. The only difference between these modes
        is the energy discretization; adding higher Legendre orders for scattering did not
        noticeably change $k$ for OpenMC-MG. The hybrid method appears to perform worse for Cases
        2b and 3b, but by the same logic, OpenMC-MG and $S_8$ appear to perform worse for Cases 2a and
        3a.

        We replaced Fig 9 with the neutron group flux distributions to illustrate the flux differences
        observed between OpenMC-CE and OpenMC-MG. We will aim to find a better multigroup structure
        for future studies.
    \end{solution}

    \question If you could get the figures closer to where they're referenced in the text, that would be most appreciated by the readers.
    \begin{solution}
        We rearranged the figures to improve proximity to their reference in the text.
    \end{solution}

    \question For many of the figures in the results section, there is extensive description of what is seen in the figures. It would be more concise to refer to the figures and give interpretations of, or meaning to, the figures.
    \begin{solution}
        We replaced extensive descriptions with concise interpretive discussion points.
    \end{solution}

    \question In Figure 8, why is the hybrid method doing better than diffusion in the diffusion regions? Shouldn't it admit the diffusion solution there? Or is it tightly coupled to the control rod subregion? Is diffusion even valid in the diffusion region? This goes back to the comment that it seems like the SN subregions should be around every interface. Conversely, why is the hybrid method not matching transport in the SN subregion?
    \begin{solution}
        The hybrid method tightly couples the $S_N$ subsolver to the diffusion solver.

        Looking closely at Fig 10, the flux shapes in the $S_N$ subregion match the transport
        solution while the flux shapes in the diffusion region matches the diffusion solution. Flux
        shapes are governed by local fission, absorption, and scattering cross sections.
        The transition from transport-like to
        diffusion-like occurs at around $x=7.5$~cm, where the buffer region starts.

        Your observations of flux mismatch point to the differing flux magnitudes due to the
        $S_N$-diffusion boundary coupling which effectively results in different neutron currents
        acting on the $S_N$ and diffusion regions at $x=7.5$~cm compared to full
        transport or full diffusion solutions.
    \end{solution}

    \question In Figure 10, do the discrepancies indicate that the SN subregion is not properly located? Or that the SN subregion should be group-dependent?
    \begin{solution}
        Drift values near the subregion boundary will always deviate from the reference values
        due to the $P_1$ approximation used for the $S_N$ boundary source from the diffusion
        solution unless the boundary is placed in a large homogeneous region appropriate for diffusion.

        With the buffer region covering $x=7.5$~cm to 10~cm, the drift values with significant
        discrepancies are not applied to the low-order diffusion model. We set up the fixed
        correction subregion to be large enough for all groups. The adaptive buffer region inner
        boundary handles group-dependent behavior.
    \end{solution}

    \question In Figure 12, what is the meaning of the non-monotonic hybrid curve?
    \begin{solution}
        The non-monotonic curves in Figure 11 and 12 indicate that keff errors from
        applying diffusion to the lattice region are spatially non-uniform. 
        As the correction subregion size increases, it encompasses more of the fuel-graphite
        lattice region where transport corrections are applied. The control rod region in Case 3b
        also induces a significant spatial shift in the flux distribution and may be contributing
        to the non-monotonic behavior.
    \end{solution}

    \question Page 33 mentioned minimizing computation cost, but there is no mention of computational cost. In Monte Carlo, a Figure of Merit of 1/[R$^2\times$ T], where R is a relative error and T is compute time, is designed to validly compare runs with different numbers of particles. However, it could be useful for you to report something like a FOM for you calculations, especially the 2D ones.
    \begin{solution}
        In response to Comment 5, we added computational runtime comparisons between the hybrid and
        neutron diffusion methods. Due to existing limitations regarding $S_N$ solver memory usage,
        we provided a baseline estimate of $S_8$ computational costs
        based on the fact that an $S_8$ simulation has 80$\times$ as many unknowns as a neutron
        diffusion simulation on the same mesh.
    \end{solution}

    \question Table III seems backward to me. If I spend more time better converging an inner iteration, I would expect the outers to converge in fewer iterations. If I converge the inners only loosely, more of the work falls upon the outers and there will be more iterations overall required. The efficiency would depend on the relative cost of inner and outer iterations. This would be a good place to include computational cost or efficiency. Can you explain?
    \begin{solution}
        The behavior in Table III is due to how outer iteration convergence is determined in
        MOOSE-based apps. In our response to Comment 15, we clarified that the outer iteration
        convergence check is performed by checking whether the residual value of the neutron diffusion
        solver falls below the fixed point tolerance value at the start of an outer iteration. If
        we tighten $S_N$ tolerance value, the $S_N$ solver attempts to resolve the drift correction
        terms more precisely, and in turn triggers more outer iterations to resolve the scalar
        flux and $k$ more precisely.

        This investigation on $S_N$ convergence threshold is inspired by a general observation by
        Adams and Larsen (2002) that transport
        correction terms depend weakly on the iterative transport solution and thus may have
        different convergence rates. Reynolds (2023) also notes this occurring in their
        quasidiffusion neutronics simulation of a 2-D axisymmetric MSRE model. The citations and
        relevant discussions are included in the revised draft.
    \end{solution}

    \question Please make the discussion on pages 36-38 much more concise. The comment, "This is likely due to statistical uncertainty..." seems a little inadequate; can you investigate it and state it more rigorously? You're normalizing to the average (clear if you move the $N_i$ to the denominator as 1/$N_i$), so you don't need all the discussion.
    \begin{solution}
        We made significant changes to the 2-D results section to replace descriptive text with
        concise interpretive text in response to Comment 21.
        We also added discussions throughout the paper addressing errors arising from energy
        discretization, core heterogeneity, and leakage.
    \end{solution}

    \question Page 39, lines 13-17, I don't understand the comment about scalar coefficients. Can you refer to your equations or somehow make this clearer and more robust?
    \begin{solution}
        We were referring to the fact that the neutron diffusion method performs poorly in
        low-scattering regions (e.g., air). The isotropic diffusion term results in a flat
        diffusive flux profile in air regions.
        We revised this paragraph to provide better interpretations of the figures.
    \end{solution}

    \question Page 43, the sentence at the very could be made much more concise or could be more interpretive rather than repeating information from a table. You say, "is effective," but is it really? Down below, you say that increasing the number of energy groups makes the problem more diffusive - how is that? This paragraph should be more rigorous. "which clarifies"... I'm not sure that it does. Can you plot the cross sections to better make this argument?
    \begin{solution}
        By ``increasing energy groups'', we were referring to lower-energy groups, not increasing
        the total number of groups. We edited this set of text to be less ambiguous and added a table
        that lists the 8-group fission cross sections to support our reasoning. We also moved the
        text to the beginning of the 2-D results \& discussion section for better flow.
    \end{solution}

\end{questions}

%---------------------------------------------------------------------
% REVIEWER 3
%---------------------------------------------------------------------
\section*{Reviewer 3}
\begin{questions}

    \question On page 6, line \# 10, ``while simultaneously retaining the computational efficiency of the neutron diffusion method''. This is not correct as this method cannot ensure the same efficiency as diffusion. The method only selectively employs transport calculations in certain areas (control rods in this case) while using diffusion in other areas. If the number of control rods increases, the computational cost should be significantly higher than that of diffusion.
    \begin{solution}
        We have removed the quoted phrase. Please also refer to the newly added Section IV.B.4 on
        computational performance where we compare the computational costs of neutron diffusion,
        $S_8$, and hybrid methods.
    \end{solution}

    \question Authors asserts in Eq. (10) that operator L* is the adjoint of L. Please check whether this saying is strict in mathematics. As I understand, if L* is an adjoint of L, Eq. 10 should not contain the boundary integrals.
    \begin{solution}
        Thank you for highlighting this error. We replaced this with an explicit definition,
        $L^*=-L_1+L_2$, and by explicitly showing the mathematical relation derived from
        integration by parts.
    \end{solution}

    \question On page 13, line \#13, Nonlinear diffusion acceleration (NDA) is applied in neutron transport problems with unstructured meshes, as reported in Computers and Mathematics with Applications 170 (2024) 142-160. This paper should be cited.
    \begin{solution}
        I added the suggested citation to the relevant section.
    \end{solution}

    \question It would be good if authors could illustrate how Eq. (30) is derived. Phi\_g,d is the angular flux of a discretized direction d, but not a function of omega\_d. I do not quite get how the P1 truncation could be used in this case. Besides, the omega term in Eq.(31) seems like a typo.
    \begin{solution}
        Thank you for catching this typo. Eq.\ 30 is supposed to be the P1 approximation for
        angular flux $\Psi_g(\vec{r},\hat{\Omega})$ along $\hat{\Omega}_d$. We edited the equation
        to correctly refer to $\Psi_g$ instead of $\Psi_{g,d}$. We use Eq.\ 30 as the basis for
        estimating $\Psi_{g,d}$ in Eq.\ 31 for the boundary condition in Eq.\ 19.
    \end{solution}

    \question Please explain why keff cannot converge when increasing the correction subregion size in Fig. 11. \& 12. Shouldn't the results become more accurate when more corrections are employed?
    \begin{solution}
        In Fig. 11 and 12, we show $k_\text{eff}$ estimates for subregion sizes ranging from 10~cm
        to 40~cm. 1-D problem domains of Case 3a and 3b are 50~cm across. We believe that neutron
        leakage error of 300~pcm on the right boundary is contributing to the discrepancy. Having
        the ``subregion''
        encompass the entire 50-cm domain would be equivalent to a full $S_8$ calculation with
        nonlinear diffusion acceleration, which is already depicted in the figures as red lines.
        Using the hybrid method for correcting neutron leakage at the outer boundaries is outside
        the scope of this paper. Instead, we plan to apply albedo boundary conditions derived from
        OpenMC/Serpent to tackle leakage errors.
    \end{solution}

    \question It would be helpful to give a schematic showing the meshes for the region, buffer region and subregion for the 2D MSRE case.
    \begin{solution}
        We made edits to Figures 5 and 6 to indicate the location of the subregion relative to the
        entire problem domain.
    \end{solution}

    \question I am curious that should the meshes for subregions decided before all the calculations are performed? And how are the sizes of subregions decided? Are they decided based on the mean free path, but of which energy group?
    \begin{solution}
        The subregion sizes were chosen based on placing the subregion boundary about one mean
        free path away from the control rod and air regions. We considered mean free paths of
        group 1 neutrons in fuel and graphite since they are the largest. We added this information
        to Sections II.C and III.B of the revised draft.
    \end{solution}

    \question Although it is seen in TABLE V that hybrid method outperforms diffusion method in terms of predicting the rod worth, but the errors are significant in keff. This seems to lower the attractiveness of the method, given the complex theory and the iterative process.
    \begin{solution}
        We agree with your assessment of the hybrid method in its current state. This is
        an exploratory work that we will iterate on.
        We are planning to tackle other sources of discrepancy in the diffusion method through
        more conventional approaches such as
        albedo boundary conditions on outer boundaries and $SP_N$ methods or other correction
        treatments in bulk regions since these existing methods fare decently in non-highly
        neutron absorbing regions. Our newly added discussions on flux distributions also indicate
        that fine angular resolution of higher energy fluxes may not be necessary for multiphysics
        simulations of thermal MSRs.

        Please also note the newly added discussions throughout the 1-D results
        section IV.A on $k_\text{eff}$ and flux errors arising from energy discretization,
        fuel-graphite heterogeneity, and leakage.
    \end{solution}

\end{questions}
\end{document}

  %
  % End of file `elsarticle-template-num.tex'.
